% =============================================================================
% 0 장 — 이 문서에 대해
% =============================================================================
\chapter{이 문서에 대해}\label{ch:about}

이 문서는 \ttHH, \HHfourb{} 탐색의 완전 강입자(fully hadronic, FH) 채널을 한곳에서 이해하기 위한 한국어 분석 노트다.
CMS 의 공식 analysis note(AN)가 아니라 분석자 본인을 위한 작업본이며, 나중에 영어 AN 을 쓸 때 뼈대가 된다.
분석의 흐름(어떤 단계가 어떤 순서로 이어지는가), 각 단계의 물리·통계적 근거(왜 그렇게 하는가), Run 2 AN 과 무엇이 같고 다른가,
코드에서 무엇이 언제 바뀌었는가, 그리고 아직 정하지 않은 것을 함께 적는다.
문서의 방향은 사용자 결정 D-2026-10-10-B 의 (4)에서 왔다\src{D-2026-10-10-B}.

\section{원칙}\label{sec:about-principles}

분석의 정의는 Run 2 의 \ttHH{} AN(SL·DL 채널)\src{AN-22-122 v26}을 따른다. FH 채널과 Run 3(2024)에 꼭 필요한 만큼만 바꾸고,
억지로 바꿀 바에는 Run 2 기준에 맞춘다\src{D-2026-10-10-B}. Run 2 AN 에 없는 FH 고유의 문제 — QCD 다중 제트 배경과 hadronic trigger —
는 \ttH(\bbbar) 분석의 FH 채널\src{AN-19-094 v20}을 선례로 삼는다. 이 원칙에서 벗어나는 곳은 모두 ``Run 2 AN 과 다른 점'' 상자에
근거와 함께 적는다.

\section{장의 구성과 관련 기록}\label{sec:about-map}

\begin{table}[htbp]\centering
\caption{장의 구성. 오른쪽 열은 그 장의 정본 기록(md 문서·코드)이다. 숫자와 상태는 그 기록에서 옮긴다.}\label{tab:about-map}
\small
\begin{tabularx}{\linewidth}{@{}l L L@{}}
\toprule
장 & 내용 & 정본 기록 \\
\midrule
\ref{ch:intro} & Higgs 자기 결합, \ttHH{} 신호, FH 말단 상태, 분석 사슬 & AN-22-122 §1–2, \texttt{README.md} \\
\ref{ch:samples} & 데이터·MC, lumi, weight, \ttbar{}+jets 분류와 stitching & \texttt{data/samples\_*.json}, \texttt{LUMI\_SOURCES.md}, \texttt{ttbarCategorization\_KR.md} \\
\ref{ch:trigger} & trigger 경로, PD 규칙, trigger SF & STEP 24·26, \texttt{TriggerStudy/} \\
\ref{ch:objects} & jet·lepton·사건 청소·PU·보정 & \texttt{EraConfig.h}, \texttt{CorrectionsManager.h}, PLAN\_v15 §9 \\
\ref{ch:btag} & b-tagging, shape SF·norm reweight, 2024 fixed-WP & STEP 14·25·27, D-2026-10-08-A, D-2026-10-10-A \\
\ref{ch:selection} & cut 순서, 영역, weight 단계, control plot & \texttt{SelectionCuts.h}, STEP 3·14·16 \\
\ref{ch:reco} & \chisq{} 짝짓기, event-shape, Tree v1, jet 정답 표지 & \texttt{TreeVars.h}, \texttt{GenMatch.h}, STEP 5·6·27 \\
\ref{ch:qcd} & QCD 다중 제트의 데이터 기반 추정 & \texttt{PLAN\_QCD\_DD.md} \\
\ref{ch:mva} & DNN·GATJA & \texttt{PLAN\_ML\_SYST.md} §9–10 \\
\ref{ch:syst} & 계통 불확도 & \texttt{PLAN\_ML\_SYST.md} §6 \\
\ref{ch:stats} & likelihood, CLs, Asimov, fit 계획 & — \\
\ref{ch:status} & 지금의 상태, 지금까지의 변경 흐름, 이번 판에서 찾은 문제 & \texttt{STATUS.md}, \texttt{CHANGELOG.md} \\
\ref{ch:plan} & 워크스트림, 순서, 결정 대기, 위험 & \texttt{ROADMAP\_FH.md}, \texttt{DECISIONS.md} \\
부록 \ref{ch:treev1} & Tree v1 의 모든 branch & \texttt{ttHHanalyzer\_unified} 의 \texttt{initTree}, \texttt{fillTreeV1\_} \\
부록 \ref{ch:glossary} & 용어 풀이 & — \\
부록 \ref{ch:sources} & 출처 목록과 표기법 & — \\
\bottomrule
\end{tabularx}
\end{table}

각 장은 대체로 같은 순서로 쓴다: (a) 물리적 배경·이론, (b) Run 2 \ttHH{} AN 에서는, (c) FH 고유의 것은 \ttH(\bbbar) FH AN 에서는,
(d) 우리 FH 분석에서는(지금의 코드), (e) Run 3(2024) 조정과 근거, (f) 변경 이력, (g) 미결 사항.

\section{표지·상자·출처 읽는 법}\label{sec:about-legend}

\paragraph{상태 표지.} 본문과 표 안의 작은 표지는 그 내용의 상태다.
\tagimpl{} 코드에 있다.\enspace
\tagtest{} 코드에 있고 시험(단위 시험·offline smoke·KNU smoke)으로 확인했다.\enspace
\tagplan{} 계획만 있다.\enspace
\tagopen{} 정하지 않았다(결정 대기).\enspace
\tagan{} Run 2 AN 의 정의를 그대로 쓴다.\enspace
\tagrunthree{} Run 3 에 맞춰 바꿨다(근거를 같이 적음).\enspace
\tagours{} AN 에 없는 우리 선택·해석이다.

\paragraph{상자.} 네 종류다.
\begin{note}[이해를 위한 설명]
유도, 직관, ``왜 이렇게 하는가''를 푼다. 건너뛰어도 흐름은 이어진다.
\end{note}
\begin{andiff}
Run 2 AN 의 정의에서 벗어난 곳과 그 근거(POG 권고, 데이터 사정, FH 의 특성).
\end{andiff}
\begin{openq}
아직 정하지 않은 것. 누가 정해야 하는지와 권고가 있으면 같이 적는다.
\end{openq}
\begin{impl}
그 내용이 구현된 파일과 함수. 코드를 찾아갈 때 쓴다.
\end{impl}

\paragraph{출처.} 숫자와 사실 뒤의 회색 괄호가 출처다. 예: $\sigma(\ttHH)=0.756\fb$\src{AN-22-122 v26, PDF p.23, Table 9}.
AN 은 PDF 쪽 번호로 적는다(인쇄 쪽 번호와 다르다). 발표는 PDF 쪽(=슬라이드 번호), 우리 기록은 결정 번호(D-…)·변경 기록(STEP n)·문서의 절,
코드는 파일과 함수 이름으로 적는다. 출처의 전체 목록과 약어는 부록~\ref{ch:sources}에 있다. 출처 괄호가 없는 문장은 교과서 수준의 설명이거나
우리의 판단이며, 판단에는 \tagours{} 또는 ``우리 판단''이라고 적었다.

\section{이 문서를 고치는 규칙}\label{sec:about-maintenance}

\begin{itemize}
\item \textbf{정본은 md 문서와 코드다.} 이 문서는 그것을 하나의 이야기로 엮는다. 숫자를 옮길 때는 항상 출처를 붙이고, 두 곳이 다르면 정본(코드·결정 기록)을 따르고 다름을 적는다.
\item \textbf{코드가 바뀌면}(새 STEP) 그 장의 ``변경 이력''과 \ref{ch:status}장의 표를 같이 고친다.
\item \textbf{결정이 생기면}(새 D-…) 그 장의 본문과 \texttt{openq} 상자, \ref{ch:plan}장의 결정 대기 목록을 고친다.
\item \textbf{판 번호}는 \texttt{main.tex} 의 \texttt{\textbackslash ANKRversion} 과 아래 개정 이력에 함께 올린다.
\item \textbf{빌드}: \file{docs/AN_KR/build.sh}(XeLaTeX, kotex, Noto CJK KR 글꼴). PDF 는 git 에 넣지 않고(\texttt{*.pdf} 는 무시됨)
맥 워크스페이스의 \file{AN_KR_pdf/} 에 둔다. 한 장만 시험할 때는 \code{build.sh --chapter sections/05_btag.tex}.
\item 작성 규칙(매크로·표지·출처 형식)은 \file{docs/AN_KR/README.md} 에 있다.
\end{itemize}

\section{이 판을 만든 방법과 한계}\label{sec:about-howmade}

v0.1 은 AI(Claude)가 원문 AN 두 편(AN-22-122 v26, AN-19-094 v20)의 본문 추출본, 발표 자료 여섯 편, 저장소의 결정·변경 기록과 코드를
읽고 초안을 썼다. 장마다 숫자에는 출처를 달았고, 출처끼리 다른 곳은 \texttt{openq} 상자에 남겼다. 그림 안의 숫자를 읽은 값에는
``그림 판독''이라고 적었다. 쓰는 동안 출처를 교차 확인하다가 분석 자체의 문제도 몇 가지 드러났다(\ref{sec:status-findings}절).
이 판은 분석자가 읽으며 고칠 초안이다. 특히 이론 설명은 교과서 수준의 요약이므로 원 논문과 함께 읽는다.

\section{개정 이력}\label{sec:about-history}

\begin{table}[htbp]\centering
\caption{개정 이력}\label{tab:about-history}
\small
\begin{tabularx}{\linewidth}{@{}l l L@{}}
\toprule
판 & 날짜 & 내용 \\
\midrule
v0.1 & 2026-10-10 & 첫 판. 0–13장과 부록 A–C. 코드 기준은 tempTTHH \texttt{dev260614} 의 STEP 27(커밋 N·O). \\
\bottomrule
\end{tabularx}
\end{table}
